\section{Introduction}

Consider two descriptions of the same world. The agents, object movements,
and false belief are identical; only the linguistic route changes. The relevant
state may appear earlier or later, under a name or an ambiguous pronoun, in an
active clause or inside an impersonal report. A language model can reverse its
answer even though the world has not changed. The usual diagnosis is prompt
brittleness. Our starting point is different: perhaps the description changed
\emph{which source was allowed to control the answer}.

Standard Theory-of-Mind evaluation cannot separate these possibilities. A
correct false-belief answer is compatible with following the queried agent,
copying reality, tracking another agent, or exploiting a cue that happens to
agree with the target. Semantics-preserving perturbations expose this fragility
\citep{ullman2023trivial,shapira2024cleverhans}, but a catalog of reversals still
identifies only the winning answer. It does not identify the controller that
made that answer win.

We treat a description as a \emph{binding controller}: a linguistic interface
that regulates access to source-indexed states while the model arbitrates among
the queried agent's belief, reality, another agent's state, and an object
distractor. This reframes the problem from ``does the model contain belief
information?'' to ``when does a description make one source govern the
answer?'' The framing connects model evaluation to discourse accessibility,
where order and referring form alter which entities or propositions remain
available for interpretation
\citep{gundel1993givenness,grosz1995centering,webster-nothman-2016-using}.

This change of object turns two failures into the main scientific clues. First,
pooled wording effects reverse across item pools. They become systematic only
after we expose the \emph{typed conflict topology}: whether the target source
competes with reality, another agent, or both. Second, one otherwise successful
linguistic contrast reverses when each event is placed under a reportative
predicate. The facts remain present, but their attributed route to the answer
changes. Together these observations suggest that linguistic form controls
source access rather than contributing a context-free answer bias.

That hypothesis makes a quantitative prediction. Passive voice, event-local
embedding, and source attribution each reduce the target source's influence.
If the reductions are independent, their costs should add in an unseen compound
construction. If they act through a shared accessibility bottleneck, the second
constraint should have less room to operate once the first route is limiting.
The latter predicts a positive compound-minus-additive residual. We freeze the
primitive costs before observing the compounds and confirm the prediction on
new latent worlds in Qwen2.5-7B-Instruct.

The law is not universal across the checkpoints we test, and that failure
sharpens rather than dissolves the result. The identical frozen protocol on
Qwen3-8B and matched Qwen2.5-7B Base retains partial control-relative predictor
advantages but not the directional residual. A positive affine change of
readout scale would preserve residual sign and normalized geometry; it does
not. The checkpoints therefore differ in the \emph{shape} of linguistic-source
composition, not merely in confidence or logit amplitude.

The paper's argument has three parts. We first introduce a typed source-control
object that makes controller identity behaviorally observable under fixed
semantics. We then derive and independently test a source-accessibility
saturation law for compound linguistic constraints. Finally, we use frozen
cross-checkpoint failures and an affine-gain invariance result to delimit the
law as checkpoint-sensitive. This supports a falsifiable behavioral theory of
language-mediated source routing. It does not identify a unique aggregation
operator, a hidden-state circuit, a causal effect of instruction tuning, or a
human psycholinguistic universal.
